\section{Experience-Dependent Structural Plasticity in Knowledge Graphs: The Bridge Twin Engine}

\textbf{Author}: Bryan Alexander Buchorn  
\textbf{Affiliation}: Independent Researcher, Las Vegas, NV, USA  
\textbf{Status}: v2.52.0 (Phase 172 (STRB) COMPLETE)
\textbf{Date}: May 2, 2026

---

\subsubsection{Abstract}
We introduce the \textbf{Bridge Twin Engine}, a mechanism for autonomous topological evolution in Knowledge Graphs (KGs) based on reasoning experience. While traditional KGs maintain a static structure, the Bridge Twin Engine implements a graph-theoretic analog to biological Long-Term Potent\-iation (LTP) \cite{hebb1949, bipoo1998}. By mater\-ializing "Twin Nodes" to act as structural relays across frequently traversed community boundaries, the system short-circuits latent paths and optimizes its own geometry for future inference. We formalize the poten\-tiation rules based on usage frequency and semantic salience, and provide a pruning mechanism (LTD) to maintain topological sanity. The v2.52.0 release incor\-porates an atomic re-mapping protocol to synchronize relays with background community re-parti\-tioning, effectively solving the "Zombie Bridge" problem in dynamic envir\-onments. As of v2.52.0, the GraphBridgeEngine (Phase 30) performs proactive cross-component bridge synthesis before queries encounter dead ends, scaling to 8,300 active bridges in WebQSP OPT confi\-guration and contr\-ibuting to H@10=20.84\% (versus 16.59\% without bridges); async bridge synthesis via TaskQueue further decouples bridge updates from active beam traversal.

\subsubsection{1. Introd\-uction}
The efficiency of multi-hop reasoning in graphs is funda\-mentally constrained by the diameter and sparsity of the network. In large-scale KGs, seman\-tically related concepts often reside in disparate community partitions, neces\-sitating high-latency "Bridge" traversals. We propose that a graph should not be a static artifact, but a plastic entity that reshapes its topology to favor successful reasoning chains.

\subsubsection{2. Methodology}

\paragraph{2.1 The Potent\-iation Rule}
A structural relay node $v'_{twin}$ is mater\-ialized for node $v$ in destination community $C_{dest}$ if:
\begin{enumerate}
\item \textbf{Frequency}: Traversal $(u \to v)$ occurs $\geq n_{min}$ times.
\item \textbf{Alignment}: $\cos(\vec{e}_v, \vec{c}_{dest}) \geq \sigma$, where $\vec{c}_{dest}$ is the community centroid.

\end{enumerate}
\paragraph{2.2 Materi\-alization and Reflection}
Upon poten\-tiation, the engine performs a "Reflection" operation:
\begin{itemize}
\item Assign $v'_{twin}$ to $C_{dest}$.
\item Establish a \texttt{BRIDGE\_TWIN} binding edge $(v, v'_{twin})$.
\item Mirror all of $v$'s edges that terminate in $C_{dest}$ onto $v'_{twin}$.

\end{itemize}
\subsubsection{3. Structural Maintenance: The LTD Analog}
To prevent the "Inform\-ational Sprawl" of mater\-ialized nodes, we implement a temporal decay function:
\begin{equation}
c_t = c_0 \cdot \lambda^{\Delta t}
\end{equation}
where $\lambda$ is the decay constant and $\Delta t$ is the time since last usage. Edges falling below a confidence threshold $c_{min}$ are pruned during the \textbf{REM Cycle} [Buchorn, 2026].

\subsubsection{4. Conclusion}
The Bridge Twin Engine provides a biolo\-gically-inspired framework for self-optimizing Knowledge Graphs. By trans\-forming reasoning history into physical graph structure, it enables sub-millisecond multi-hop reasoning on incre\-asingly complex network topologies. In CEREBRUM v2.52.0, proactive bridge synthesis scales to 8,300 active bridges in the WebQSP OPT confi\-guration, contr\-ibuting directly to a H@10 improvement from 16.59\% to 20.84\% — a concrete, quantified demon\-stration of experience-dependent structural plasticity improving reasoning recall.

---

\subsection{5. Recent Advances (v2.52.0 -\textgreater  v2.52.0)}

The Bridge Twin Engine has been subst\-antially extended since v2.51.1. The following describes the most significant devel\-opments.

\textbf{GraphBridgeEngine: Proactive Cross-Component Synthesis (Phase 30).} The original Bridge Twin Engine was reactive — bridges mater\-ialized only after a traversal encountered a community boundary. The \texttt{GraphBridgeEngine} (Phase 30) introduces proactive synthesis: at graph load time and after each community re-parti\-tioning, the engine identifies disco\-nnected or weakly connected components and pre-mater\-ializes bridge nodes across their boundaries. This eliminates the cold-start penalty where early queries must "earn" bridges before benefiting from them.

\textbf{Scale: 8,300 Active Bridges on WebQSP.} In the WebQSP OPT confi\-guration, the system maintains appro\-ximately 8,300 active bridge records. The impact on recall is measurable: H@10 rises from 16.59\% (without bridge synthesis) to 20.84\% (with GraphBridgeEngine), a relative gain of +25.6\%.

\textbf{Async Bridge Synthesis via TaskQueue (Phase 39).} Prior to Phase 39, bridge mater\-ialization occurred synch\-ronously in the beam traversal hot path, adding latency to queries that triggered new poten\-tiation events. Phase 39 decouples this via a \texttt{TaskQueue}: poten\-tiation events are enqueued and processed by a background worker. Active traversals are unaffected by bridge write operations, and the community map swap post-rebalance remains atomic.

\textbf{Post-Rebalance Bridge Pruning (Phase 19).} After the \texttt{GlobalRebalancer} triggers a background DSCF/TSC re-run and performs the atomic community map swap, a post-rebalance hook notifies \texttt{BridgeTwinEngine} to prune stale bridge records whose source or destination community assignments have changed. This prevents "Zombie Bridges" accum\-ulating after topology shifts.

\textbf{Benchmark Impact.}

\begin{center}
\begin{tabularx}{\columnwidth}{|>{\raggedright\arraybackslash}X|>{\raggedright\arraybackslash}X|}
\hline
Config\-uration & H@10 (WebQSP) \\ \hline
No bridge synthesis & 16.59\% \\ \hline
With GraphBridgeEngine (OPT) & 20.84\% \\ \hline
Gain & +25.6\% relative \\ \hline
\end{tabularx}
\end{center}

---
\subsection{Acknow\-ledgments}

The author gratefully ackno\-wledges the use of Claude (Anthropic) as a research assistant throughout this work. Claude assisted with mathe\-matical forma\-lization, code generation, manuscript preparation, and technical writing. All conceptual contr\-ibutions, archi\-tectural decisions, exper\-imental design, and intel\-lectual claims are solely the author's.

\textbf{References}
\begin{enumerate}
\item Hebb, D. O. (1949). The organ\-ization of behavior: A neuro\-psychological theory.
\item Bi, G. Q., \& Poo, M. M. (1998). Synaptic modif\-ications in cultured hippocampal neurons. Journal of Neuros\-cience.
\item Markram, H., et al. (1997). Regulation of synaptic efficacy by predictions of spike timing. Science.
\item Newman, M. E. (2006). Modularity and community structure in networks. PNAS.
\item Abbott, L. F., \& Nelson, S. B. (2000). Synaptic plasticity: taming the beast. Nature Neuros\-cience.
\item Caporale, N., \& Dan, Y. (2008). Spike timing-dependent plasticity: a Hebbian learning rule. Annual Review of Neuros\-cience.
\item Buchorn, B. A. (2026). CEREBRUM v2.51: Complete Technical Specif\-ication for Autonomous Knowledge Graph Reasoning. [CEREBRUM\_REPORT\_PLACEHOLDER].

\end{enumerate}
---
\textbf{Reviewed on}: May 2, 2026 for version v2.52.0
